% =========================================================
% Closure Rules for Families of Sets
% =========================================================

\subsection{Closure Rules for Families of Sets}

\begin{exposition}
A set operation may exist without a chosen family being closed under that
operation.  If \(A,B\subseteq X\), then \(A\cup B\) exists as a subset of
\(X\).  A separate question is whether \(A\cup B\) belongs to a particular
family \(\mathcal F\subseteq\mathcal P(X)\).
\end{exposition}

\begin{definitionbox}{Definition (Closed Under Set Operations)}
\begin{definition}[Closed Under Set Operations]\label{def:closed-under-set-operation-vocabulary}
Let \(X\) be a set and let \(\mathcal F\subseteq\mathcal P(X)\).
\begin{enumerate}[label=(\roman*)]
\item \(\mathcal F\) is \emph{closed under a unary operation} \(T\) if
      \(A\in\mathcal F\) implies \(T(A)\in\mathcal F\).
\item \(\mathcal F\) is \emph{closed under a binary operation} \(\star\) if
      \(A,B\in\mathcal F\) implies \(A\star B\in\mathcal F\).
\item \(\mathcal F\) is \emph{closed under an indexed operation} if every
      indexed family \((A_i)_{i\in I}\) of members of \(\mathcal F\) has its
      indexed output again in \(\mathcal F\).
\item \(\mathcal F\) is \emph{closed under a collection operation} if every
      subcollection \(\mathcal G\subseteq\mathcal F\) has its collection output
      again in \(\mathcal F\).
\end{enumerate}
\end{definition}
\end{definitionbox}

\begin{remark*}[Standard quantified statement]
\[
\begin{aligned}
\mathcal F\text{ is closed under }T
&\iff \forall A\in\mathcal F,\ T(A)\in\mathcal F,\\
\mathcal F\text{ is closed under }\star
&\iff \forall A,B\in\mathcal F,\ A\star B\in\mathcal F,\\
\mathcal F\text{ is closed under an indexed operation}
&\iff \forall (A_i)_{i\in I}\subseteq\mathcal F,\ \mathcal O((A_i)_{i\in I})\in\mathcal F.
\end{aligned}
\]
\end{remark*}

\begin{remark*}[Predicate reading]
\[
\begin{aligned}
\operatorname{ClosedUnderUnaryOperation}(\mathcal F,T)
&\iff \forall A\in\mathcal F,\ T(A)\in\mathcal F,\\
\operatorname{ClosedUnderBinaryOperation}(\mathcal F,\star)
&\iff \forall A,B\in\mathcal F,\ A\star B\in\mathcal F.
\end{aligned}
\]
\end{remark*}

\begin{remark*}[Interpretation]
Closure is a stability property.  The operation produces a subset of \(X\);
closure says the family admits that output as one of its selected members.
\end{remark*}

\begin{remark*}[Exposition]
The following phrases use the preceding definition:
\[
\begin{array}{ll}
\text{closed under complements}
  & A\in\mathcal F\Rightarrow X\setminus A\in\mathcal F,\\
\text{closed under pairwise unions}
  & A,B\in\mathcal F\Rightarrow A\cup B\in\mathcal F,\\
\text{closed under pairwise intersections}
  & A,B\in\mathcal F\Rightarrow A\cap B\in\mathcal F,\\
\text{closed under countable unions}
  & A_n\in\mathcal F\ \forall n\in\mathbb N
    \Rightarrow \bigcup_{n\in\mathbb N}A_n\in\mathcal F,\\
\text{closed under arbitrary unions}
  & \mathcal G\subseteq\mathcal F
    \Rightarrow \bigcup_{G\in\mathcal G}G\in\mathcal F.
\end{array}
\]
\end{remark*}

\LeanFormalizes{def:closed-under-set-operation-vocabulary}{lra-lean}{LRA.VolumeI.Set.Families}{ClosedUnderUnaryOperation}{pending}
\LeanFormalizes{def:closed-under-set-operation-vocabulary}{lra-lean}{LRA.VolumeI.Set.Families}{ClosedUnderBinaryOperation}{pending}
\LeanFormalizes{def:closed-under-set-operation-vocabulary}{lra-lean}{LRA.VolumeI.Set.Families}{ClosedUnderIndexedOperation}{pending}
\LeanFormalizes{def:closed-under-set-operation-vocabulary}{lra-lean}{LRA.VolumeI.Set.Families}{ClosedUnderCollectionOperation}{pending}
\LeanFormalizes{def:closed-under-set-operation-vocabulary}{lra-lean}{LRA.VolumeI.Set.Families}{ClosedUnderComplements}{pending}
\LeanFormalizes{def:closed-under-set-operation-vocabulary}{lra-lean}{LRA.VolumeI.Set.Families}{ClosedUnderPairwiseUnions}{pending}
\LeanFormalizes{def:closed-under-set-operation-vocabulary}{lra-lean}{LRA.VolumeI.Set.Families}{ClosedUnderPairwiseIntersections}{pending}
\LeanFormalizes{def:closed-under-set-operation-vocabulary}{lra-lean}{LRA.VolumeI.Set.Families}{ClosedUnderPairwiseDifferences}{pending}
\LeanFormalizes{def:closed-under-set-operation-vocabulary}{lra-lean}{LRA.VolumeI.Set.Families}{ClosedUnderPairwiseSymmetricDifferences}{pending}
\LeanFormalizes{def:closed-under-set-operation-vocabulary}{lra-lean}{LRA.VolumeI.Set.Families}{ClosedUnderCountableUnions}{pending}
\LeanFormalizes{def:closed-under-set-operation-vocabulary}{lra-lean}{LRA.VolumeI.Set.Families}{ClosedUnderCountableIntersections}{pending}
\LeanFormalizes{def:closed-under-set-operation-vocabulary}{lra-lean}{LRA.VolumeI.Set.Families}{ClosedUnderArbitraryUnions}{pending}
\LeanFormalizes{def:closed-under-set-operation-vocabulary}{lra-lean}{LRA.VolumeI.Set.Families}{ClosedUnderArbitraryIntersections}{pending}
\LeanFormalizes{def:closed-under-set-operation-vocabulary}{lra-lean}{LRA.VolumeI.Set.Families}{FiniteUnionFromList}{pending}
\LeanFormalizes{def:closed-under-set-operation-vocabulary}{lra-lean}{LRA.VolumeI.Set.Families}{FiniteIntersectionFromList}{pending}
\LeanFormalizes{def:closed-under-set-operation-vocabulary}{lra-lean}{LRA.VolumeI.Set.Families}{ClosedUnderFiniteUnions}{pending}
\LeanFormalizes{def:closed-under-set-operation-vocabulary}{lra-lean}{LRA.VolumeI.Set.Families}{ClosedUnderFiniteIntersections}{pending}

\begin{dependencies}
\begin{itemize}
  \item \hyperref[def:indexed-family]{Indexed Family of Sets}
  \item \hyperref[def:union]{Union}
  \item \hyperref[def:intersection]{Intersection}
  \item \hyperref[def:complement]{Complement}
  \item \hyperref[def:subset]{Subset}
\end{itemize}
\end{dependencies}

\begin{theorembox}{Theorem (Closure Rule Intersection Stability)}
\begin{theorem}[Closure Rule Intersection Stability]\label{thm:closure-rule-intersection-stability}
Let \(\mathfrak S\subseteq\mathcal P(\mathcal P(X))\) be a collection system.
If every \(\mathcal F\in\mathfrak S\) is closed under a fixed unary, binary,
indexed, or collection operation, then
\[
\bigcap_{\mathcal F\in\mathfrak S}\mathcal F
\]
is closed under the same operation.

In particular, intersection preserves closure under complements, pairwise
unions, pairwise intersections, countable unions, countable intersections,
arbitrary unions, and arbitrary intersections.
\hyperref[prf:closure-rule-intersection-stability]{Go to proof.}
\end{theorem}
\end{theorembox}

\begin{remark*}[Standard quantified statement]
\[
\forall \mathcal T\subseteq\mathfrak S,\quad
\bigl(\forall \mathcal F\in\mathcal T,\ \mathcal F\text{ is closed under }\mathcal O\bigr)
\to
\bigcap_{\mathcal F\in\mathcal T}\mathcal F\text{ is closed under }\mathcal O.
\]
\end{remark*}

\begin{remark*}[Predicate reading]
\[
\forall \mathcal T\subseteq\mathfrak S,\quad
\bigl(\forall \mathcal F\in\mathcal T,\ \mathcal F\text{ is closed under }\mathcal O\bigr)
\to
\bigcap_{\mathcal F\in\mathcal T}\mathcal F\text{ is closed under }\mathcal O.
\]
The predicate-level content is stability of the same closure rule under
intersection of the eligible families.
\end{remark*}

\begin{remark*}[Interpretation]
This is the mechanism behind generated structures.  If all admissible
structures obey a closure rule, then the intersection of all admissible
structures also obeys that rule.
\end{remark*}

\LeanFormalizes{thm:closure-rule-intersection-stability}{lra-lean}{LRA.VolumeI.Set.Families}{UnaryClosureStableUnderCollectionIntersection}{pending}
\LeanFormalizes{thm:closure-rule-intersection-stability}{lra-lean}{LRA.VolumeI.Set.Families}{BinaryClosureStableUnderCollectionIntersection}{pending}
\LeanFormalizes{thm:closure-rule-intersection-stability}{lra-lean}{LRA.VolumeI.Set.Families}{IndexedClosureStableUnderCollectionIntersection}{pending}
\LeanFormalizes{thm:closure-rule-intersection-stability}{lra-lean}{LRA.VolumeI.Set.Families}{CollectionClosureStableUnderCollectionIntersection}{pending}

\begin{dependencies}
\begin{itemize}
  \item \hyperref[def:closed-under-set-operation-vocabulary]{Closed Under Set Operations}
  \item \hyperref[def:collection-system]{Collection System}
  \item \hyperref[def:indexed-intersection]{Indexed Intersection}
  \item \hyperref[def:subset]{Subset}
\end{itemize}
\end{dependencies}

\begin{theorembox}{Theorem (Pairwise Closure Implies Finite Closure)}
\begin{theorem}[Pairwise Closure Implies Finite Closure]\label{thm:pairwise-closure-implies-finite-closure}
Let \(\mathcal F\subseteq\mathcal P(X)\).
\begin{enumerate}[label=(\roman*)]
\item If \(\varnothing\in\mathcal F\) and \(\mathcal F\) is closed under
      pairwise unions, then \(\mathcal F\) is closed under finite unions.
\item If \(X\in\mathcal F\) and \(\mathcal F\) is closed under pairwise
      intersections, then \(\mathcal F\) is closed under finite intersections.
\end{enumerate}
\hyperref[prf:pairwise-closure-implies-finite-closure]{Go to proof.}
\end{theorem}
\end{theorembox}

\begin{remark*}[Standard quantified statement]
\[
\begin{aligned}
\varnothing\in\mathcal F\ \wedge\
\mathcal F\text{ is closed under pairwise unions}
&\to
\mathcal F\text{ is closed under finite unions},\\
X\in\mathcal F\ \wedge\
\mathcal F\text{ is closed under pairwise intersections}
&\to
\mathcal F\text{ is closed under finite intersections}.
\end{aligned}
\]
\end{remark*}

\begin{remark*}[Interpretation]
Pairwise closure is the two-input closure rule.  Finite closure is obtained by
iterating the two-input rule over a finite list.
\end{remark*}

\begin{remark*}[Exposition]
The empty finite union is \(\varnothing\), and the empty finite intersection is
\(X\).  These conventions explain why the theorem includes
\(\varnothing\in\mathcal F\) for finite unions and \(X\in\mathcal F\) for
finite intersections.
\end{remark*}

\LeanFormalizes{thm:pairwise-closure-implies-finite-closure}{lra-lean}{LRA.VolumeI.Set.Families}{PairwiseUnionClosureImpliesFiniteUnionClosure}{pending}
\LeanFormalizes{thm:pairwise-closure-implies-finite-closure}{lra-lean}{LRA.VolumeI.Set.Families}{PairwiseIntersectionClosureImpliesFiniteIntersectionClosure}{pending}

\begin{dependencies}
\begin{itemize}
  \item \hyperref[def:closed-under-set-operation-vocabulary]{Closed Under Set Operations}
  \item \hyperref[def:empty-set]{Empty Set}
  \item \hyperref[def:union]{Union}
  \item \hyperref[def:intersection]{Intersection}
\end{itemize}
\end{dependencies}

\begin{theorembox}{Theorem (Complement-Union-Intersection Closure Duality)}
\begin{theorem}[Complement-Union-Intersection Closure Duality]\label{thm:complement-union-intersection-closure-duality}
Let \(\mathcal F\subseteq\mathcal P(X)\) be closed under complements.
\begin{enumerate}[label=(\roman*)]
\item If \(\mathcal F\) is closed under pairwise unions, then
      \(\mathcal F\) is closed under pairwise intersections.
\item If \(\mathcal F\) is closed under pairwise intersections, then
      \(\mathcal F\) is closed under pairwise unions.
\item If \(\mathcal F\) is closed under finite unions, then
      \(\mathcal F\) is closed under finite intersections.
\item If \(\mathcal F\) is closed under finite intersections, then
      \(\mathcal F\) is closed under finite unions.
\end{enumerate}
\hyperref[prf:complement-union-intersection-closure-duality]{Go to proof.}
\end{theorem}
\end{theorembox}

\begin{remark*}[Standard quantified statement]
\[
\begin{aligned}
\mathcal F\text{ complement-closed and union-closed}
&\to \mathcal F\text{ intersection-closed},\\
\mathcal F\text{ complement-closed and intersection-closed}
&\to \mathcal F\text{ union-closed},
\end{aligned}
\]
at the corresponding pairwise or finite level.
\end{remark*}

\begin{remark*}[Interpretation]
Complement closure and De Morgan laws make union closure and intersection
closure interchangeable at the same finite level.
\end{remark*}

\LeanFormalizes{thm:complement-union-intersection-closure-duality}{lra-lean}{LRA.VolumeI.Set.Families}{ComplementAndPairwiseUnionClosureImpliesPairwiseIntersectionClosure}{pending}
\LeanFormalizes{thm:complement-union-intersection-closure-duality}{lra-lean}{LRA.VolumeI.Set.Families}{ComplementAndPairwiseIntersectionClosureImpliesPairwiseUnionClosure}{pending}
\LeanFormalizes{thm:complement-union-intersection-closure-duality}{lra-lean}{LRA.VolumeI.Set.Families}{ComplementAndFiniteUnionClosureImpliesFiniteIntersectionClosure}{pending}
\LeanFormalizes{thm:complement-union-intersection-closure-duality}{lra-lean}{LRA.VolumeI.Set.Families}{ComplementAndFiniteIntersectionClosureImpliesFiniteUnionClosure}{pending}

\begin{dependencies}
\begin{itemize}
  \item \hyperref[def:closed-under-set-operation-vocabulary]{Closed Under Set Operations}
  \item \hyperref[thm:de-morgan]{De Morgan Laws}
  \item \hyperref[thm:involution]{Involution}
  \item \hyperref[thm:pairwise-closure-implies-finite-closure]{Pairwise Closure Implies Finite Closure}
\end{itemize}
\end{dependencies}
